\question[4]
Gegeben ist der folgende Python-Code:

\begin{lstlisting}
summe = 0
for i in range(10):
    summe = summe + 0.1
print(summe == 1.0)
\end{lstlisting}

\begin{parts}
\part Das Programm gibt \texttt{False} aus. Erläutere, warum.
\part Ändere die letzte Zeile so, dass geprüft wird, ob \texttt{summe} ungefähr gleich 1.0 ist.
\end{parts}

\begin{solutionbox}{5cm}
\begin{parts}
\part Gleitkommazahlen werden binär mit begrenzter Genauigkeit gespeichert. 0.1 lässt sich im Binärsystem nicht exakt
darstellen. Bei jeder Addition entsteht ein kleiner Rundungsfehler, am Ende ist \texttt{summe} gleich
\texttt{0.9999999999999999} und nicht genau 1.0.
\part
\begin{lstlisting}
print(abs(summe - 1.0) < 0.000001)
\end{lstlisting}
oder mit \texttt{import math}: \texttt{print(math.isclose(summe, 1.0))}
\end{parts}
\end{solutionbox}
